% =====================================================================
\chapter{Soluciones}
\label{ap:soluciones}
% =====================================================================

Las soluciones de la parte~B se presentan apartado por apartado, en cinco partes: datos, planteamiento, desarrollo, resultado y comprobación. Se indican también los criterios de corrección cuando el reparto de puntos de un apartado no es evidente. Todos los valores numéricos se han comprobado con Python; en los pasos intermedios se conservan cuatro cifras significativas.

% ---------------------------------------------------------------------
\section{Clave de respuestas de la parte A}

\begin{table}[H]
\centering\footnotesize
\renewcommand{\arraystretch}{1.05}
\caption{Clave de respuestas de la parte~A.}\label{tab:clave}
\begin{tabularx}{\textwidth}{@{}clX>{\raggedright\arraybackslash}p{1.95cm}@{}}
\toprule
N.º & Resp. & Justificación & Apuntes \\
\midrule
A.1 & (c) & BPE cuenta los pares adyacentes, ponderados por frecuencia, y fusiona el más frecuente. & M12, §2.1 \\
A.2 & (b) & Parte de los 256 valores posibles de un byte; por eso nunca aparece un token desconocido. & M12, §2.4 \\
A.3 & (c) & «p», «i» y «n» (ASCII) ocupan 1~byte cada una; «ñ» y «ó», 2 cada una: $3\cdot 1+2\cdot 2=7$. & M12, §1.3 \\
A.4 & (a) & Las reglas se aplican en el orden en que se aprendieron; así la codificación es determinista. & M12, §2.2 \\
A.5 & (d) & La recurrencia $h_t=f(h_{t-1},x_t)$ exige $n$ pasos secuenciales y el gradiente se multiplica $n$ veces. & M12, §5.5 \\
A.6 & (c) & Cada token (fila) tiene un peso para cada token (columna). & M13, §2.1 \\
A.7 & (b) & $q\cdot k$ es la suma de $d_k$ productos independientes de media~0 y varianza~1. & M13, §2.3 \\
A.8 & (d) & $e^{-\infty}=0$ y los pesos permitidos siguen sumando~1; un~0 antes de la softmax no anula nada. & M13, §3.1 \\
A.9 & (d) & $4d^{2}=4\cdot 768^{2}=2\,359\,296$, sin dependencia de $h$; (c) es $2d^{2}$. & M13, §4.2 \\
A.10 & (b) & La memoria de $A$ crece como $n^{2}$: $(8192/1024)^{2}=8^{2}=64$. & M13, §4.4 \\
A.11 & (a) & Pre-norm normaliza la entrada de la subcapa y deja limpio el camino residual; (b) es post-norm. & M14, §2.2 \\
A.12 & (a) & En pre-norm, $h^{L}=h^{0}+\sum$ aportaciones y su norma crece con $L$. & M14, §2.3 \\
A.13 & (b) & $3d\,\dff=2d\cdot 4d=8d^{2}$, luego $\dff=\tfrac{8}{3}d$. & M14, §3.3 \\
A.14 & (a) & $R_m\T R_n=R_{n-m}$, luego $\tilde q_m\cdot\tilde k_n=q\T R_{n-m}k$. & M14, §4.3 \\
A.15 & (c) & $v$ es el contenido que se copia; la posición solo debe decidir los pesos. & M14, §4.3 \\
A.16 & (c) & $2Vd=2\cdot 32\,000\cdot 4096=262\,144\,000$. & M14, §1.4 y §5.2 \\
A.17 & (d) & $2+2+4+4+4=16$~bytes por parámetro. & M14, §5.3 \\
A.18 & (d) & $6ND=6\cdot 10^{9}\cdot 2\cdot 10^{10}=1{,}2\cdot 10^{20}$. & M14, §6.3 \\
A.19 & (a) & Logits iguales dan la distribución uniforme y la pérdida $\ln V=\ln 50\,257=10{,}82$. & M15, §1.4 \\
A.20 & (b) & La pérdida de validación estima el rendimiento sobre datos nuevos. & M15, §4.2 y §4.4 \\
\bottomrule
\end{tabularx}
\end{table}

% ---------------------------------------------------------------------
\section{Soluciones de la parte B}

% ===================================================================== B1
\begin{solucion}{eje:B-bloque}

\apartado{a}{5}
\datos Transformer decoder-only del estilo de LLaMA: $n$ tokens, dimensión $d$, $L$ bloques, $h$ cabezas, SwiGLU con $\dff$, RMSNorm, RoPE y $V$ tokens sin atar.
\planteamiento recorrer el camino de los datos: enteros, embedding, $L$ bloques pre-norm, normalización final, unembedding y softmax (M14, §1.2).
\pasos
\begin{enumerate}[label=\arabic*.]
  \item Entrada: los $n$ enteros seleccionan filas de $W_E\in\R^{V\times d}$ y forman $h^{0}\in\R^{n\times d}$.
  \item Primera subcapa del bloque: RMSNorm~1, proyecciones $Q$, $K$, $V$, rotación RoPE de $q$ y $k$ (aquí entra la posición), atención causal en cada cabeza, concatenación y $W_O$, y suma residual~1.
  \item Segunda subcapa: RMSNorm~2, SwiGLU con $W_{\mathrm{gate}}$, $W_{\mathrm{up}}$ y $W_{\mathrm{down}}$, y suma residual~2.
  \item Tras los $L$ bloques: RMSNorm final, $W_U\in\R^{d\times V}$ y softmax.
\end{enumerate}
\resultado \textbf{el diagrama de la \cref{fig:bloque}}, con dos residuales por bloque, una normalización antes de cada subcapa, RoPE aplicada a $q$ y $k$ y una RMSNorm final.
\comprobacion todas las flechas entre bloques transportan matrices $n\times d$; solo cambian de forma la entrada (enteros a $\R^{d}$) y la salida ($\R^{d}$ a $\R^{V}$).

\begin{figure}[ht]
\centering
\begin{tikzpicture}[
  caja/.style={draw,rounded corners=2pt,minimum width=7.2cm,minimum height=0.62cm,font=\small,align=center,fill=white},
  suma/.style={draw,circle,inner sep=0.5pt,font=\small,fill=white},
  flecha/.style={-{Stealth[length=2.2mm]}},
  node distance=0.36cm]
\node[caja] (tok) {Tokens $x_1,\dots,x_n$ (enteros)};
\node[caja,below=of tok] (emb) {Embedding $W_E$ ($V\times d$) $\;\to\; h^{0}$ ($n\times d$)};
\node[caja,below=0.75cm of emb] (n1) {RMSNorm 1};
\node[caja,below=of n1] (qkv) {$Q=XW_Q$, $K=XW_K$, $V=XW_V$};
\node[caja,below=of qkv] (rope) {RoPE sobre $q$ y $k$ (entra la posición)};
\node[caja,below=of rope] (att) {$h$ cabezas: $\softmax\bigl(Q_iK_i\T/\sqrt{d_k}+M\bigr)V_i$};
\node[caja,below=of att] (wo) {Concatenar y $W_O$ ($d\times d$)};
\node[suma,below=of wo] (s1) {$+$};
\node[caja,below=of s1] (n2) {RMSNorm 2};
\node[caja,below=of n2] (mlp) {SwiGLU: $\bigl(\SiLU(XW_{\mathrm{gate}})\odot XW_{\mathrm{up}}\bigr)W_{\mathrm{down}}$};
\node[suma,below=of mlp] (s2) {$+$};
\node[caja,below=0.75cm of s2] (nf) {RMSNorm final};
\node[caja,below=of nf] (wu) {Unembedding $W_U$ ($d\times V$): logits $n\times V$};
\node[caja,below=of wu] (sm) {softmax $\to P(x_{n+1}\mid x_{1:n})$};
\foreach \a/\b in {tok/emb,emb/n1,n1/qkv,qkv/rope,rope/att,att/wo,wo/s1,s1/n2,n2/mlp,mlp/s2,s2/nf,nf/wu,wu/sm}
  \draw[flecha] (\a) -- (\b);
\coordinate (r0) at ($(emb.south)!0.5!(n1.north)$);
\draw[flecha] (r0) -- ++(-4.3,0) |- (s1.west) node[pos=0.25,left,font=\footnotesize,align=right] {residual\\1};
\coordinate (r1) at ($(s1.south)!0.5!(n2.north)$);
\draw[flecha] (r1) -- ++(-4.3,0) |- (s2.west) node[pos=0.25,left,font=\footnotesize,align=right] {residual\\2};
\node[draw,dashed,rounded corners=4pt,inner sep=5pt,fit=(n1)(qkv)(rope)(att)(wo)(s1)(n2)(mlp)(s2)] (bl) {};
\node[anchor=north west,font=\footnotesize\itshape] at (bl.north east) {$\times L$};
\end{tikzpicture}
\caption{Transformer decoder-only con bloques pre-norm del estilo de LLaMA. El recuadro discontinuo es un bloque y se repite $L$ veces.}\label{fig:bloque}
\end{figure}

Criterios: embedding y unembedding con sus formas (1~p); dos residuales que puentean cada subcapa (1~p); normalizaciones antes de cada subcapa y normalización final (1~p); atención causal con máscara y RoPE en $q$ y $k$ (1~p); MLP con sus tres matrices (1~p). Si se dibujan posiciones aprendidas sumadas al embedding, al estilo de GPT-2, se acepta siempre que se indique.

\apartado{b}{3}
\datos entrada del bloque $h^{\ell-1}$; arquitectura pre-norm.
\planteamiento en pre-norm, cada subcapa lee su entrada normalizada y suma su salida al flujo residual (M14, §2.2).
\pasos
\begin{enumerate}[label=\arabic*.]
  \item Atención: $a^{\ell}=h^{\ell-1}+\MHA\bigl(\RMSNorm_1(h^{\ell-1})\bigr)$.
  \item MLP: $h^{\ell}=a^{\ell}+\SwiGLU\bigl(\RMSNorm_2(a^{\ell})\bigr)$.
  \item Salida: $\mathrm{logits}=\RMSNorm_f(h^{L})\,W_U$.
\end{enumerate}
\resultado \textbf{las tres ecuaciones anteriores}, 1~p cada una.
\comprobacion desarrollando la recurrencia, $h^{L}=h^{0}+\sum_{\ell}\bigl(\text{aportación de la atención}+\text{aportación del MLP}\bigr)$: el camino identidad no atraviesa ninguna normalización (M14, §2.1).

\apartado{c}{3}
\datos $n$ tokens, $d$, $h$ cabezas, $\dff$, $V$.
\planteamiento seguir las formas en cada producto $Y=XW$; 0,5~p por forma.
\pasos
\begin{enumerate}[label=\arabic*.]
  \item Entrada al primer bloque: $n$ filas de $W_E$, es decir, $n\times d$.
  \item $Q=XW_Q$ es $n\times d$; separada en $h$ cabezas de dimensión $d/h$: $h\times n\times (d/h)$.
  \item Atención de una cabeza: $(n\times d/h)(d/h\times n)$, es decir, $n\times n$.
  \item $W_{\mathrm{gate}}\in\R^{d\times\dff}$: $(n\times d)(d\times\dff)$ da $n\times\dff$.
  \item Salida de un bloque: $n\times d$, la misma forma que la entrada.
  \item Logits: $(n\times d)(d\times V)$ da $n\times V$.
\end{enumerate}
\resultado \textbf{$n\times d$; $h\times n\times(d/h)$; $n\times n$; $n\times\dff$; $n\times d$; $n\times V$.}
\comprobacion como $h\cdot(d/h)=d$, la concatenación de las cabezas recupera la forma $n\times d$.

\apartado{d}{2}
\datos subcapas de atención y MLP.
\planteamiento comprobar si la operación combina filas distintas de la matriz $n\times d$ (M14, §3.1).
\pasos
\begin{enumerate}[label=\arabic*.]
  \item La atención mezcla tokens: la fila $i$ de $O=AV$ es $\sum_{j\le i}a_{ij}v_j$.
  \item El MLP aplica la misma función a cada fila por separado: la salida del token~$i$ solo depende de la fila~$i$.
  \item Cada subcapa \emph{suma} su salida al flujo residual, que vive en $\R^{d}$; la suma exige la misma dimensión, y por eso $W_O\in\R^{d\times d}$ y $W_{\mathrm{down}}\in\R^{\dff\times d}$ devuelven a $\R^{d}$.
\end{enumerate}
\resultado \textbf{la atención mezcla información entre tokens; el MLP actúa token a token; ambas devuelven $\R^{d}$ porque se suman al flujo residual.}
\comprobacion si $W_{\mathrm{down}}$ fuera $\dff\times\dff$, la suma $a^{\ell}+\SwiGLU(\cdot)$ no estaría definida.

\apartado{e}{2}
\datos pre-norm, $h=x+F(\Norm(x))$, frente a post-norm, $h=\Norm(x+F(x))$.
\planteamiento comparar el camino del gradiente y la estabilidad del entrenamiento (M14, §2.2).
\pasos
\begin{enumerate}[label=\arabic*.]
  \item En pre-norm existe un camino identidad de la entrada a la salida: el jacobiano $\partial h^{L}/\partial h^{0}$ contiene el término $I$ y el gradiente llega a las primeras capas sin atravesar normalizaciones.
  \item El entrenamiento es estable con muchas capas y con un calentamiento corto o nulo; en post-norm, el gradiente de las últimas capas es grande en la inicialización y obliga a un calentamiento largo y cuidadoso~\cite{xiong}.
\end{enumerate}
\resultado \textbf{gradiente limpio por el camino identidad y entrenamiento estable en profundidad, poco sensible al calentamiento}, 1~p por razón.
\comprobacion GPT-2~\cite{gpt2}, LLaMA-2~\cite{llama2} y la tarea~1 de CS336~\cite[§3.4]{cs336} usan pre-norm; el problema \emph{pre\_norm\_ablation} de esa tarea compara ambas variantes experimentalmente.

\repasar{M14, §1.2, §2.1, §2.2, §2.3 y §3.1}
\end{solucion}

% ===================================================================== B2
\begin{solucion}{eje:B-atencion}

\apartado{a}{3}
\datos $X$ con filas $x_1=(1,0)$, $x_2=(0,1)$, $x_3=(1,1)$; $W_Q,W_K\in\R^{2\times 4}$, $W_V\in\R^{2\times 2}$.
\planteamiento la fila $i$ de $XW$ es $x_iW$: $x_1$ selecciona la primera fila de $W$, $x_2$ la segunda y $x_3$ suma ambas (M13, §2.1).
\pasos
\begin{enumerate}[label=\arabic*.]
  \item $Q$: $(1,1,0,0)$, $(0,0,1,1)$ y su suma $(1,1,1,1)$.
  \item $K$: $(2,0,0,0)$, $(0,0,2,0)$ y su suma $(2,0,2,0)$.
  \item $V$: $(1,1)$, $(0,1)$ y su suma $(1,2)$.
\end{enumerate}
\resultado
\[
Q=\begin{pmatrix}1&1&0&0\\0&0&1&1\\1&1&1&1\end{pmatrix},\qquad
K=\begin{pmatrix}2&0&0&0\\0&0&2&0\\2&0&2&0\end{pmatrix},\qquad
V=\begin{pmatrix}1&1\\0&1\\1&2\end{pmatrix},
\]
1~p cada matriz.
\comprobacion en las tres matrices, la fila~3 es la suma de las filas 1 y~2, porque $x_3=x_1+x_2$.

\apartado{b}{3}
\datos $Q$ y $K$ del apartado~(a); $d_k=4$.
\planteamiento $S_{ij}=q_i\cdot k_j$; después se divide por $\sqrt{d_k}=2$.
\pasos
\begin{enumerate}[label=\arabic*.]
  \item Fila 1, $q_1=(1,1,0,0)$: $q_1\cdot k_1=2$, $q_1\cdot k_2=0$, $q_1\cdot k_3=2$.
  \item Fila 2, $q_2=(0,0,1,1)$: $q_2\cdot k_1=0$, $q_2\cdot k_2=2$, $q_2\cdot k_3=2$.
  \item Fila 3, $q_3=(1,1,1,1)$: $q_3\cdot k_1=2$, $q_3\cdot k_2=2$, $q_3\cdot k_3=4$.
\end{enumerate}
\resultado
\[
S=\begin{pmatrix}2&0&2\\0&2&2\\2&2&4\end{pmatrix},\qquad
\frac{S}{\sqrt{d_k}}=\begin{pmatrix}1&0&1\\0&1&1\\1&1&2\end{pmatrix};
\]
2~p por $S$ y 1~p por el escalado.
\comprobacion la columna~3 de $S$ es la suma de las columnas 1 y~2, porque $k_3=k_1+k_2$.

\apartado{c}{2}
\datos tres posiciones; máscara causal.
\planteamiento $M_{ij}=0$ si $j\le i$ y $M_{ij}=-\infty$ si $j>i$ (M13, §3.1).
\pasos
\begin{enumerate}[label=\arabic*.]
  \item $M$ es triangular superior estricta de $-\infty$; se suma elemento a elemento a $S/\sqrt{d_k}$:
  \[
  M=\begin{pmatrix}0&-\infty&-\infty\\0&0&-\infty\\0&0&0\end{pmatrix},\qquad
  Z=\begin{pmatrix}1&-\infty&-\infty\\0&1&-\infty\\1&1&2\end{pmatrix}.
  \]
  \item Con $-\infty$, $e^{-\infty}=0$: el peso de cada token futuro es exactamente nulo y los pesos permitidos se renormalizan solos. Con~$0$, $e^{0}=1$ y el futuro conservaría peso; por ejemplo, la fila~1 sería $\softmax(1,0,1)=(0{,}422;\allowbreak\ 0{,}155;\allowbreak\ 0{,}422)$.
\end{enumerate}
\resultado \textbf{$M$ y $Z$ como arriba; $-\infty$ anula exactamente los pesos futuros}, mientras que $0$ no anula nada.
\comprobacion la fila~$i$ de $Z$ conserva $i$ valores finitos: 1, 2 y~3.

\apartado{d}{3}
\datos $Z$ del apartado~(c); $e\approx 2{,}718$, $e^{2}\approx 7{,}389$; $V$ del apartado~(a).
\planteamiento softmax por filas con $e^{-\infty}=0$; después $O=AV$.
\pasos
\begin{enumerate}[label=\arabic*.]
  \item Fila 1: solo $z_{11}$ es finito, luego $a_1=(1,0,0)$.
  \item Fila 2: $e^{0}=1$ y $e^{1}=2{,}718$, con suma $3{,}718$: $a_{21}=0{,}269$ y $a_{22}=0{,}731$.
  \item Fila 3: dividiendo $e^{1},e^{1},e^{2}$ por $e$ quedan $1,1,e$, con suma $4{,}718$: $a_{31}=a_{32}=0{,}212$ y $a_{33}=0{,}576$.
  \item $o_1=(1,1)$; $o_2=0{,}269\,(1,1)+0{,}731\,(0,1)=(0{,}269;\allowbreak\ 1)$.
  \item $o_3=0{,}212\,(1,1)+0{,}212\,(0,1)+0{,}576\,(1,2)=(0{,}788;\allowbreak\ 1{,}576)$, con los pesos a cuatro cifras ($0{,}2119$ y $0{,}5761$).
\end{enumerate}
\resultado
\[
A=\begin{pmatrix}1&0&0\\0{,}269&0{,}731&0\\0{,}212&0{,}212&0{,}576\end{pmatrix},\qquad
O=AV=\begin{pmatrix}1&1\\0{,}269&1\\0{,}788&1{,}576\end{pmatrix}.
\]
\comprobacion cada fila de $A$ suma~1; como la segunda componente de $v_1$ y de $v_2$ es~1, la de $o_2$ vale exactamente~1. Criterios: máscara aplicada antes de la softmax (1~p), $A$ correcta (1~p), $O$ correcta (1~p).

\apartado{e}{2}
\datos $x_3=(5,-3)$; el resto, igual.
\planteamiento $A$ es triangular inferior: $o_1$ y $o_2$ solo dependen de los tokens 1 y~2 (M13, §3.2). Solo hay que recalcular la fila~3.
\pasos
\begin{enumerate}[label=\arabic*.]
  \item $o_1$ y $o_2$ no dependen de $x_3$: solo cambia $o_3$.
  \item $q_3=5\,(1,1,0,0)-3\,(0,0,1,1)=(5,5,-3,-3)$; $k_3=5\,(2,0,0,0)-3\,(0,0,2,0)=(10,0,-6,0)$; $v_3=5\,(1,1)-3\,(0,1)=(5,2)$.
  \item Puntuaciones escaladas: $q_3\cdot k_1/2=10/2=5$; $q_3\cdot k_2/2=-6/2=-3$; $q_3\cdot k_3/2=(50+18)/2=34$.
  \item $e^{34}/e^{5}=e^{29}\approx 3{,}9\cdot 10^{12}$, luego $a_{33}\approx 1$, $a_{31}\approx e^{-29}\approx 2{,}5\cdot 10^{-13}$ y $a_{32}\approx e^{-37}\approx 8{,}5\cdot 10^{-17}$.
  \item $o_3\approx v_3=(5,2)$.
\end{enumerate}
\resultado \textbf{solo cambia $o_3$, que pasa a ser $(5{,}000;\allowbreak\ 2{,}000)$}: la cabeza \emph{copia} el valor $v_3$, porque la puntuación de $k_3$ es mucho mayor que las demás.
\comprobacion con Python, $1-a_{33}=2{,}5\cdot 10^{-13}$. Es la situación de «copia en la atención» de la tarea~3 de CS224N~\cite{cs224n}: casi todo el peso recae en un solo elemento cuando $q\cdot k_j\gg q\cdot k_i$ para todo $i\ne j$.

\apartado{f}{2}
\datos $a_{22}=0{,}731$, del apartado~(d).
\planteamiento para la softmax, $\partial a_{ij}/\partial z_{ij}=a_{ij}(1-a_{ij})$; el escalado por $\sqrt{d_k}$ evita la saturación~\cite[§3.2.1]{vaswani} (M13, §2.3).
\pasos
\begin{enumerate}[label=\arabic*.]
  \item $\partial a_{22}/\partial z_{22}=0{,}731\cdot 0{,}269=0{,}197$.
  \item Si las puntuaciones crecen sin escalar, la softmax se satura: $a_{ij}\to 0$ o $a_{ij}\to 1$, y en ambos casos $a_{ij}(1-a_{ij})\to 0$.
  \item Sin gradiente, la cabeza deja de aprender a quién atender.
\end{enumerate}
\resultado \textbf{$\partial a_{22}/\partial z_{22}=0{,}197$; con puntuaciones grandes tiende a~0 (softmax saturada).}
\comprobacion el máximo de $a(1-a)$ es $0{,}25$, en $a=0{,}5$, y $0{,}197<0{,}25$. Respecto de la puntuación sin escalar $s=2z$, la regla de la cadena añade el factor $1/\sqrt{d_k}=1/2$: $0{,}098$.

\repasar{M13, §2.1, §2.3, §3.1 y §3.2}
\end{solucion}

% ===================================================================== B3
\begin{solucion}{eje:B-bpe}

\apartado{a}{1}
\datos \tok{pino}~$\times 4$, \tok{pinos}~$\times 2$, \tok{vino}~$\times 3$, \tok{vinos}~$\times 1$, \tok{pinta}~$\times 2$.
\planteamiento número de símbolos $=\sum(\text{longitud}\times\text{frecuencia})$.
\pasos
\begin{enumerate}[label=\arabic*.]
  \item $4\cdot 4+5\cdot 2+4\cdot 3+5\cdot 1+5\cdot 2=16+10+12+5+10$.
\end{enumerate}
\resultado \textbf{53 símbolos.}
\comprobacion hay $4+2+3+1+2=12$ palabras de 4 o 5 letras: el total debe estar entre 48 y~60.

\apartado{b}{6}
\datos el corpus del apartado~(a), partido en caracteres.
\planteamiento en cada paso se cuentan los pares adyacentes ponderados por la frecuencia de la palabra, se fusiona el más frecuente y se restan sus apariciones del total de símbolos~\cite{sennrich} (M12, §2.1). 1,5~p por paso.
\pasos
\begin{enumerate}[label=\arabic*.]
  \item Recuento inicial: $(\tok{i},\tok{n})=4+2+3+1+2=12$; $(\tok{n},\tok{o})=4+2+3+1=10$; $(\tok{p},\tok{i})=4+2+2=8$; $(\tok{v},\tok{i})=4$; $(\tok{o},\tok{s})=3$. Gana $(\tok{i},\tok{n})$; símbolos: $53-12=41$.
  \item Tras \tok{in}: $(\tok{in},\tok{o})=10$; $(\tok{p},\tok{in})=8$; $(\tok{v},\tok{in})=4$; $(\tok{o},\tok{s})=3$. Gana $(\tok{in},\tok{o})$; símbolos: $41-10=31$.
  \item Tras \tok{ino}: $(\tok{p},\tok{ino})=4+2=6$; $(\tok{v},\tok{ino})=4$; $(\tok{ino},\tok{s})=3$. Gana $(\tok{p},\tok{ino})$; símbolos: $31-6=25$.
  \item Tras \tok{pino}: $(\tok{v},\tok{ino})=4$; empatan con~2 $(\tok{pino},\tok{s})$, $(\tok{p},\tok{in})$, $(\tok{in},\tok{t})$ y $(\tok{t},\tok{a})$. Gana $(\tok{v},\tok{ino})$; símbolos: $25-4=21$.
\end{enumerate}

\begin{center}\small
\begin{tabularx}{\textwidth}{@{}cll>{\raggedright\arraybackslash}Xc@{}}
\toprule
Paso & Par (frec.) & Segundo (frec.) & Corpus tras la fusión & Símbolos \\
\midrule
1 & \tok{i n} (12) & \tok{n o} (10) & \tok{p in o}$\times4$, \tok{p in o s}$\times2$, \tok{v in o}$\times3$, \tok{v in o s}$\times1$, \tok{p in t a}$\times2$ & 41 \\
2 & \tok{in o} (10) & \tok{p in} (8) & \tok{p ino}$\times4$, \tok{p ino s}$\times2$, \tok{v ino}$\times3$, \tok{v ino s}$\times1$, \tok{p in t a}$\times2$ & 31 \\
3 & \tok{p ino} (6) & \tok{v ino} (4) & \tok{pino}$\times4$, \tok{pino s}$\times2$, \tok{v ino}$\times3$, \tok{v ino s}$\times1$, \tok{p in t a}$\times2$ & 25 \\
4 & \tok{v ino} (4) & cuatro pares (2) & \tok{pino}$\times4$, \tok{pino s}$\times2$, \tok{vino}$\times3$, \tok{vino s}$\times1$, \tok{p in t a}$\times2$ & 21 \\
\bottomrule
\end{tabularx}
\end{center}

\resultado \textbf{reglas, en orden: \tok{in}, \tok{ino}, \tok{pino}, \tok{vino}; el corpus queda con 21 símbolos.}
\comprobacion recuento directo: $1\cdot 4+2\cdot 2+1\cdot 3+2\cdot 1+4\cdot 2=21$. No hay empates en el primer puesto, de modo que la regla de desempate (M12, §2.3) no interviene. Error frecuente: tomar $(\tok{p},\tok{i})=8$ como máximo olvidando las apariciones en \tok{vino}.

\apartado{c}{3}
\datos reglas en orden: $(\tok{i},\tok{n})$, $(\tok{in},\tok{o})$, $(\tok{p},\tok{ino})$, $(\tok{v},\tok{ino})$.
\planteamiento partir cada palabra en caracteres y aplicar las reglas en el orden aprendido (M12, §2.2).
\pasos
\begin{enumerate}[label=\arabic*.]
  \item \tok{vinos}: \tok{v i n o s} $\to$ \tok{v in o s} $\to$ \tok{v ino s} $\to$ \tok{vino s}: 2 tokens.
  \item \tok{pinto}: \tok{p i n t o} $\to$ \tok{p in t o}; la regla~2 no se aplica (tras \tok{in} viene \tok{t}) y la~3 necesita \tok{ino}: 4 tokens.
  \item \tok{sino}: \tok{s i n o} $\to$ \tok{s in o} $\to$ \tok{s ino}; no existe la regla $(\tok{s},\tok{ino})$: 2 tokens.
  \item \tok{ovino}: \tok{o v i n o} $\to$ \tok{o v in o} $\to$ \tok{o v ino} $\to$ \tok{o vino}: 2 tokens.
\end{enumerate}
\resultado \textbf{\tok{vinos} $=$ [\tok{vino}, \tok{s}] (2); \tok{pinto} $=$ [\tok{p}, \tok{in}, \tok{t}, \tok{o}] (4); \tok{sino} $=$ [\tok{s}, \tok{ino}] (2); \tok{ovino} $=$ [\tok{o}, \tok{vino}] (2).}
\comprobacion ninguna palabra nueva produce un token desconocido, porque los caracteres sueltos pertenecen siempre al vocabulario.

\apartado{d}{2}
\datos «piñón» en UTF-8.
\planteamiento los caracteres ASCII ocupan 1~byte; «ñ» y «ó», 2 (M12, §1.3 y §1.5).
\pasos
\begin{enumerate}[label=\arabic*.]
  \item Bytes en hexadecimal: \texttt{70}, \texttt{69}, \texttt{C3\,B1}, \texttt{C3\,B3}, \texttt{6E}: $1+1+2+2+1=7$.
  \item En español, los pares $(\texttt{C3},\texttt{B1})$, que forman «ñ», y $(\texttt{C3},\texttt{B3})$, que forma «ó» (y análogos como $(\texttt{C3},\texttt{A1})$ para «á»), son muy frecuentes.
  \item Sin esas fusiones, ninguna palabra con tilde o con «ñ» puede llegar a ser un único token.
\end{enumerate}
\resultado \textbf{7 símbolos iniciales; conviene fusionar pronto $(\tok{C3},\tok{B1})$ y $(\tok{C3},\tok{B3})$}, los dos bytes de «ñ» y de «ó».
\comprobacion coincide con la cuestión~\ref{q:pinon} (7~bytes).

\repasar{M12, §1.3, §1.5, §2.1, §2.2 y §2.4}
\end{solucion}

% ===================================================================== B4
\begin{solucion}{eje:B-lince}

\apartado{a}{2}
\datos $d=2048$.
\planteamiento $\dff=\tfrac{8}{3}d$ redondeado hacia arriba al múltiplo de 256~\cite{shazeer} (M14, §3.3).
\pasos
\begin{enumerate}[label=\arabic*.]
  \item $8\cdot 2048/3=16\,384/3=5461{,}33$.
  \item $5461{,}33/256=21{,}33$; el entero siguiente es~22.
  \item $22\cdot 256=5632$.
\end{enumerate}
\resultado \textbf{$\dff=5632$.}
\comprobacion $21\cdot 256=5376<5461{,}33<5632$.

\apartado{b}{4}
\datos $d=2048$, $L=24$, $\dff=5632$, $V=48\,000$; embeddings atados; RMSNorm solo con pesos; sin sesgos.
\planteamiento $N=L\,(4d^{2}+3d\,\dff+2d)+Vd+d$, con una sola matriz de embeddings por estar atados (M14, §5.1 y §5.2).
\pasos
\begin{enumerate}[label=\arabic*.]
  \item $d^{2}=4\,194\,304$; atención: $4d^{2}=16\,777\,216$.
  \item MLP: $3\cdot 2048\cdot 5632=34\,603\,008$.
  \item Dos RMSNorm: $2\cdot 2048=4096$. Bloque: $16\,777\,216+34\,603\,008+4096=51\,384\,320$.
  \item 24 bloques: $24\cdot 51\,384\,320=1\,233\,223\,680$.
  \item Embedding: $48\,000\cdot 2048=98\,304\,000$; RMSNorm final: $2048$.
  \item $N=1\,233\,223\,680+98\,304\,000+2048=1\,331\,529\,728$.
  \item $12Ld^{2}=12\cdot 24\cdot 4\,194\,304=1\,207\,959\,552$, que es el $90{,}7\,\%$ de $N$.
  \item MLP: $24\cdot 34\,603\,008=830\,472\,192$, el $62{,}4\,\%$ de $N$.
\end{enumerate}

\begin{center}\small
\begin{tabular}{@{}llr@{}}
\toprule
Pieza & Fórmula & Parámetros \\
\midrule
Atención & $4d^{2}$ & $16\,777\,216$ \\
MLP SwiGLU & $3d\,\dff$ & $34\,603\,008$ \\
Dos RMSNorm & $2d$ & $4096$ \\
\midrule
\textbf{Bloque} & & $\mathbf{51\,384\,320}$ \\
\bottomrule
\end{tabular}
\end{center}

\resultado \textbf{bloque: $51\,384\,320$; $N=1\,331\,529\,728\approx 1{,}33\cdot 10^{9}$; $12Ld^{2}$ es el $90{,}7\,\%$ de $N$; los MLP suponen el $62{,}4\,\%$.}
\comprobacion la diferencia con la aproximación procede del embedding ($98\,304\,000/N=7{,}4\,\%$) y del redondeo de $\dff$ hacia arriba. Sin atar, $N$ sería $1\,429\,833\,728$.

\apartado{c}{3}
\datos $N=1{,}3315\cdot 10^{9}$; $D=2\cdot 10^{11}$; H100 a $989\cdot 10^{12}$~FLOP/s; aprovechamiento del 40\,\%; 64~GPU.
\planteamiento inferencia, $\approx 2N$ FLOPs por token; entrenamiento, $C\approx 6ND$~\cite{kaplan}; tiempo $=C/\text{velocidad efectiva}$ (M14, §6.2 y §6.3).
\pasos
\begin{enumerate}[label=\arabic*.]
  \item Inferencia: $2\cdot 1{,}3315\cdot 10^{9}=2{,}663\cdot 10^{9}$ FLOPs por token.
  \item Entrenamiento: $C=6\cdot 1{,}3315\cdot 10^{9}\cdot 2\cdot 10^{11}=1{,}598\cdot 10^{21}$ FLOPs.
  \item Velocidad efectiva por GPU: $0{,}4\cdot 989\cdot 10^{12}=3{,}956\cdot 10^{14}$ FLOP/s.
  \item Tiempo: $1{,}598\cdot 10^{21}/3{,}956\cdot 10^{14}=4{,}039\cdot 10^{6}$~s, es decir, $1122$~h de GPU.
  \item Con 64~GPU: $1122/64=17{,}53$~h de reloj ($0{,}73$~días).
\end{enumerate}
\resultado \textbf{$2{,}66\cdot 10^{9}$ FLOPs por token; $C\approx 1{,}60\cdot 10^{21}$ FLOPs; unas $1122$~h de H100; $17{,}5$~h de reloj con 64~GPU.}
\comprobacion LLaMA-2~7B~\cite{llama2} ($N=6{,}74\cdot 10^{9}$, $D=2\cdot 10^{12}$) cuesta $8{,}09\cdot 10^{22}$ FLOPs, unas $50{,}6$ veces más, coherente con un $N$ cinco veces mayor y un $D$ diez veces mayor.

\apartado{d}{3}
\datos $N=1\,331\,529\,728$; bf16, 2~bytes; int4, $0{,}5$~bytes; Adam en precisión mixta, 16~bytes por parámetro.
\planteamiento memoria $=N\times$ bytes por parámetro (M14, §5.3).
\pasos
\begin{enumerate}[label=\arabic*.]
  \item bf16: $2N=2{,}663\cdot 10^{9}$~bytes $=2{,}66$~GB.
  \item int4: $0{,}5N=6{,}658\cdot 10^{8}$~bytes $=0{,}67$~GB.
  \item Adam en precisión mixta: $16N=2{,}130\cdot 10^{10}$~bytes $=21{,}3$~GB.
  \item $21{,}3<24$: los estados caben, pero solo quedan unos $2{,}7$~GB para las activaciones.
\end{enumerate}
\resultado \textbf{$2{,}66$~GB en bf16, $0{,}67$~GB en int4 y $21{,}3$~GB para el estado del entrenamiento.} El estado cabe en 24~GB, pero las activaciones, que crecen con el lote y el contexto, no caben con holgura.
\comprobacion $21{,}3/2{,}66=8=16/2$. En la práctica se recurriría a la recomputación de activaciones o al reparto entre varias GPU (ZeRO, módulo~17).

\apartado{e}{3}
\datos $L=24$, $d=2048$, $n=4096$; $2N=2{,}663\cdot 10^{9}$.
\planteamiento comparar $4Lnd$ con $2N$ y despejar $n$ en la igualdad: $n^{*}=N/(2Ld)$ (M14, §6.2).
\pasos
\begin{enumerate}[label=\arabic*.]
  \item $4\cdot 24\cdot 4096\cdot 2048=805\,306\,368\approx 8{,}053\cdot 10^{8}$ FLOPs por token.
  \item Fracción: $8{,}053\cdot 10^{8}/2{,}663\cdot 10^{9}=0{,}302$.
  \item Igualdad: $n^{*}=2N/(4Ld)=2{,}663\cdot 10^{9}/196\,608=13\,545$.
\end{enumerate}
\resultado \textbf{la atención añade un $30{,}2\,\%$ a $2N$ con $n=4096$, y lo iguala con $n\approx 13\,545$ tokens.}
\comprobacion la fracción es lineal en $n$: $4096/13\,545=0{,}302$.

\repasar{M14, §3.3, §5.1, §5.2, §5.3, §6.2 y §6.3}
\end{solucion}

% ===================================================================== B5
\begin{solucion}{eje:B-rope}

\apartado{a}{5}
\datos $d_k=2$, $\theta=\pi/4$, $q=(1,0)$, $k=(1,1)$; $\cos(\pi/4)=\sin(\pi/4)=0{,}707$.
\planteamiento calcular las dos constantes de la identidad y evaluar en cada $\varphi=(n-m)\theta$~\cite{roformer} (M14, §4.3).
\pasos
\begin{enumerate}[label=\arabic*.]
  \item $q\cdot k=1$ y $q_1k_2-q_2k_1=1$; la puntuación es $\cos\varphi-\sin\varphi$.
  \item $(2,1)$ y $(5,4)$: $n-m=-1$, $\varphi=-\pi/4$, puntuación $0{,}707+0{,}707=1{,}414$.
  \item $(3,3)$: $\varphi=0$, puntuación $1$.
  \item $(1,2)$: $\varphi=\pi/4$, puntuación $0$.
  \item $(0,4)$: $\varphi=\pi$, puntuación $-1$.
\end{enumerate}

\begin{center}\small
\begin{tabular}{@{}cccc@{}}
\toprule
$(m,n)$ & $n-m$ & $\varphi$ & $\cos\varphi-\sin\varphi$ \\
\midrule
$(2,1)$ & $-1$ & $-\pi/4$ & $1{,}414$ \\
$(5,4)$ & $-1$ & $-\pi/4$ & $1{,}414$ \\
$(3,3)$ & $0$ & $0$ & $1$ \\
$(1,2)$ & $+1$ & $\pi/4$ & $0$ \\
$(0,4)$ & $+4$ & $\pi$ & $-1$ \\
\bottomrule
\end{tabular}
\end{center}

\resultado \textbf{$1{,}414$; $1{,}414$; $1$; $0$; $-1$.} Los dos primeros pares ocupan posiciones distintas a la misma distancia y dan la misma puntuación: RoPE codifica la posición relativa.
\comprobacion con distancia nula se recupera $q\cdot k=1$; además, $\cos\varphi-\sin\varphi=\sqrt2\cos(\varphi+\pi/4)$, cuyo máximo, $\sqrt2=1{,}414$, se alcanza en $\varphi=-\pi/4$. Criterios: 1~p por las dos primeras filas, 1~p por cada una de las otras tres; la observación es necesaria para obtener la puntuación completa.

\apartado{b}{2}
\datos $d=2048$, $h=16$.
\planteamiento $d_k=d/h$; los parámetros de la atención son $4d^{2}$, sin dependencia de $h$ (M13, §4.2).
\pasos
\begin{enumerate}[label=\arabic*.]
  \item $d_k=2048/16=128$.
  \item Con $h=32$, $d_k=64$, pero las 32 cabezas siguen formando $W_Q$, $W_K$, $W_V$ y $W_O$ de tamaño $d\times d$: $4d^{2}=16\,777\,216$ en ambos casos.
\end{enumerate}
\resultado \textbf{$d_k=128$; duplicar $h$ no añade parámetros.}
\comprobacion coincide con la cuestión~\ref{q:mha}: el recuento $4d^{2}$ no depende de~$h$.

\apartado{c}{4}
\datos $n=32\,768$; fp16, 2~bytes; $h=16$; $L=24$; $d=2048$.
\planteamiento cada matriz $A$ es $n\times n$; FLOPs cuadráticos $4n^{2}d$ frente a lineales $8nd^{2}$ (M13, §4.3 y §4.4).
\pasos
\begin{enumerate}[label=\arabic*.]
  \item Una cabeza: $n^{2}=1\,073\,741\,824$ elementos; por 2~bytes, $2{,}147\cdot 10^{9}$~bytes $=2{,}15$~GB.
  \item Todas las cabezas y capas: $2{,}147\cdot 10^{9}\cdot 16\cdot 24=8{,}246\cdot 10^{11}$~bytes $\approx 825$~GB.
  \item Cruce: $4n^{2}d>8nd^{2}$ equivale, dividiendo por $4nd>0$, a $n>2d=4096$.
  \item FlashAttention calcula la softmax por bloques y nunca almacena $A$ completa: la memoria adicional es lineal en $n$, aunque el número de FLOPs sigue siendo $O(n^{2}d)$~\cite{dao}.
\end{enumerate}
\resultado \textbf{$2{,}15$~GB por cabeza y unos 825~GB en total; la atención domina en FLOPs a partir de $n=4096$; FlashAttention evita almacenar~$A$.}
\comprobacion 825~GB supera con mucho los 80~GB de una GPU actual: sin FlashAttention, ese contexto es inviable.

\repasar{M13, §4.2, §4.3 y §4.4; M14, §4.3}
\end{solucion}

% ===================================================================== B6
\begin{solucion}{eje:B-minigpt}

\apartado{a}{3}
\datos $V=80$, $T=128$, $d=128$, $L=4$; MLP de $4d$; LayerNorm con peso y sesgo; posiciones aprendidas; embeddings atados; lineales sin sesgo.
\planteamiento bloque: $4d^{2}+8d^{2}+2\cdot 2d$; fuera de los bloques: $Vd$ (una vez), $Td$ y la LayerNorm final, $2d$ (M15, §2.3).
\pasos
\begin{enumerate}[label=\arabic*.]
  \item $d^{2}=16\,384$; atención: $65\,536$; MLP: $8d^{2}=131\,072$.
  \item Dos LayerNorm: $2\cdot 2\cdot 128=512$. Bloque: $197\,120$.
  \item Cuatro bloques: $788\,480$.
  \item Embedding: $80\cdot 128=10\,240$; posiciones: $128\cdot 128=16\,384$; LayerNorm final: $256$.
  \item $N=788\,480+10\,240+16\,384+256=815\,360$.
\end{enumerate}
\resultado \textbf{$N=815\,360$ parámetros.}
\comprobacion $12Ld^{2}=786\,432$, el $96{,}5\,\%$ de $N$.

\apartado{b}{2}
\datos $V=80$; pérdida inicial $\approx 4{,}39$.
\planteamiento logits casi iguales dan una softmax casi uniforme y una pérdida $\ln V$ (M15, §1.4).
\pasos
\begin{enumerate}[label=\arabic*.]
  \item Con pesos pequeños, los logits son casi iguales y $p_i\approx 1/80$.
  \item $-\ln(1/80)=\ln 80=\ln 8+\ln 10=2{,}079+2{,}303=4{,}382$.
\end{enumerate}
\resultado \textbf{$\ln 80=4{,}382\approx 4{,}39$.}
\comprobacion una pérdida inicial mucho mayor indicaría una inicialización defectuosa, por ejemplo embeddings $\mathcal N(0,1)$ atados.

\apartado{c}{3}
\datos \cref{tab:curvas}; bigramas, $2{,}45$; $\ln 2=0{,}693$.
\planteamiento se elige el mínimo de validación; $\PPL=e^{\ell}$ y bits por carácter $=\ell/\ln 2$ (M15, §1.3 y §4.2).
\pasos
\begin{enumerate}[label=\arabic*.]
  \item Mínimo de validación: paso 2500, con $1{,}44$ (antes, $1{,}46$; después, $1{,}45$).
  \item $\PPL=e^{1{,}44}=4{,}221$.
  \item Bits por carácter: $1{,}44/0{,}693=2{,}078$.
  \item Bigramas: $e^{2{,}45}=11{,}59$; cociente $11{,}59/4{,}221=2{,}75$.
\end{enumerate}
\resultado \textbf{punto de control del paso 2500; $\PPL\approx 4{,}22$ y $2{,}08$ bits por carácter; el modelo de bigramas tiene $\PPL\approx 11{,}59$}: el mini-GPT duda entre casi tres veces menos opciones.
\comprobacion $2^{2{,}078}=4{,}22$: la perplejidad no depende de la base del logaritmo.

\apartado{d}{2}
\datos $B=32$, $T=128$, 5000~pasos; corpus de $10^{6}$ caracteres.
\planteamiento tokens $=B\cdot T\cdot\text{pasos}$; épocas $=$ tokens $/$ tamaño del corpus (M15, §1.2 y §4.4).
\pasos
\begin{enumerate}[label=\arabic*.]
  \item $32\cdot 128=4096$ tokens por paso; por 5000 pasos, $20\,480\,000$ tokens.
  \item Épocas: $20\,480\,000/10^{6}=20{,}48$.
  \item Brecha entre validación y entrenamiento: $1{,}44-1{,}30=0{,}14$ en el paso 2500 y $1{,}62-0{,}98=0{,}64$ en el 5000.
  \item El modelo ve cada carácter unas 20 veces y memoriza: la validación sube desde el paso 2500.
\end{enumerate}
\resultado \textbf{$20\,480\,000$ tokens, unas $20{,}5$ épocas: sobreajuste.} Remedios: más datos (el mejor), parada temprana, \emph{dropout}, más decaimiento de pesos o un modelo más pequeño.
\comprobacion la pérdida de entrenamiento baja mientras la de validación sube: es la firma del sobreajuste en la tabla de síntomas de M15, §4.4.

\apartado{e}{2}
\datos logits $(2,1,0)$; $\tau=1$ y $\tau=0{,}5$.
\planteamiento $p=\softmax(z/\tau)$ (M15, §5.2).
\pasos
\begin{enumerate}[label=\arabic*.]
  \item $\tau=1$: $e^{2}=7{,}389$, $e^{1}=2{,}718$, $e^{0}=1$, con suma $11{,}107$: $p=(0{,}665;\allowbreak\ 0{,}245;\allowbreak\ 0{,}090)$.
  \item $\tau=0{,}5$: $z/\tau=(4,2,0)$; $e^{4}=54{,}60$, $e^{2}=7{,}389$, $e^{0}=1$, con suma $62{,}99$: $p=(0{,}867;\allowbreak\ 0{,}117;\allowbreak\ 0{,}016)$.
  \item Bajar $\tau$ concentra la masa en el carácter más probable: las muestras son más correctas, pero más repetitivas. Subir $\tau$ (con $\tau=2$, $p=(0{,}506;\allowbreak\ 0{,}307;\allowbreak\ 0{,}186)$) da más variedad y más errores.
\end{enumerate}
\resultado \textbf{$\tau=1$: $(0{,}665;\allowbreak\ 0{,}245;\allowbreak\ 0{,}090)$; $\tau=0{,}5$: $(0{,}867;\allowbreak\ 0{,}117;\allowbreak\ 0{,}016)$.} Bajar la temperatura hace las muestras más predecibles y repetitivas.
\comprobacion ambas distribuciones suman~1 y, cuando $\tau\to 0^{+}$, la distribución tiende a $(1,0,0)$.

\repasar{M15, §1.2, §1.3, §1.4, §2.3, §4.2, §4.4 y §5.2}
\end{solucion}
